\chapter*{PHỤ LỤC A: CÁC KHỐI MÃ NGUỒN CỐT LÕI}
\addcontentsline{toc}{chapter}{Phụ lục A: Các khối mã nguồn cốt lõi}
\setcounter{section}{0}
\renewcommand{\thesection}{A\arabic{section}}

Phụ lục này trích dẫn nguyên bản các đoạn mã nguồn đảm nhận nhiệm vụ then chốt trong việc duy trì kiến trúc bảo mật Zero-Knowledge của hệ thống Secret Letter. Toàn bộ mã nguồn đầy đủ được công bố tại \url{https://github.com/vmcchooky/secret-letter}.

% --- PHẦN A1 ---
\section*{A1. Khởi tạo và mã hóa bằng AES-256-GCM (Web Crypto API)}
\addcontentsline{toc}{section}{A1. Khởi tạo và mã hóa bằng AES-256-GCM (Web Crypto API)}

Trích xuất từ tệp \texttt{frontend/web-app/src/lib/crypto.ts}. Đoạn mã thể hiện quá trình sinh khóa ngẫu nhiên 256-bit (cho phép trích xuất thành Raw key để ghép vào URL), tạo Nonce ngẫu nhiên và mã hóa dữ liệu văn bản.

\begin{minted}[language=typescript, caption={Hàm sinh khóa, Nonce và mã hóa AES-GCM}, label={lst:crypto-encrypt}]{typescript}
const ALGORITHM = "AES-GCM";
const KEY_LENGTH = 256;
const NONCE_LENGTH = 12; // 12 bytes = 96 bits

export async function generateKey(): Promise<CryptoKey> {
  return await crypto.subtle.generateKey(
    { name: ALGORITHM, length: KEY_LENGTH },
    true, // extractable để lấy raw bytes đưa lên URL Fragment
    ["encrypt", "decrypt"]
  );
}

export function generateNonce(): Uint8Array {
  return crypto.getRandomValues(new Uint8Array(NONCE_LENGTH));
}

export async function encryptSecret(
  plaintext: string,
  key: CryptoKey,
  nonce: Uint8Array
): Promise<Uint8Array> {
  const encoder = new TextEncoder();
  const plaintextBytes = encoder.encode(plaintext);

  const ciphertextBuffer = await crypto.subtle.encrypt(
    { name: ALGORITHM, iv: asBufferSource(nonce) },
    key,
    asBufferSource(plaintextBytes)
  );
  return new Uint8Array(ciphertextBuffer);
}
\end{minted}

% --- PHẦN A2 ---
\section*{A2. Giải mã bằng thuật toán AES-256-GCM}
\addcontentsline{toc}{section}{A2. Giải mã bằng thuật toán AES-256-GCM}

Trích xuất từ \texttt{frontend/web-app/src/lib/crypto.ts}. Hàm nhận đầu vào là đoạn bản mã (Ciphertext) lấy từ Backend và Khóa giải mã cùng Nonce bóc tách từ URL Fragment.

\begin{minted}[language=typescript, caption={Hàm giải mã AES-GCM tại trình duyệt}, label={lst:crypto-decrypt}]{typescript}
export async function decryptSecret(
  ciphertext: Uint8Array,
  key: CryptoKey,
  nonce: Uint8Array
): Promise<string> {
  const plaintextBuffer = await crypto.subtle.decrypt(
    { name: ALGORITHM, iv: asBufferSource(nonce) },
    key,
    asBufferSource(ciphertext)
  );

  const decoder = new TextDecoder();
  return decoder.decode(plaintextBuffer);
}
\end{minted}

% --- PHẦN A3 ---
\section*{A3. Kịch bản Lua nguyên tử (Atomic) bảo vệ chống Tương tranh}
\addcontentsline{toc}{section}{A3. Kịch bản Lua nguyên tử (Atomic) bảo vệ chống Tương tranh}

Trích xuất từ tệp \texttt{backend/internal/secret/redis\_service.go}. Kịch bản này được thực thi đơn luồng trên Redis, giúp hệ thống thay đổi trạng thái (từ \texttt{active} sang \texttt{opening}) một cách an toàn mà không bị vướng hiện tượng Race Condition khi nhiều yêu cầu truy xuất cùng lúc.

\begin{minted}[language=lua, caption={Khối Lua Script xác thực và cấp quyền truy xuất (claimOpenSecretScript)}, label={lst:redis-lua}]{lua}
local metaRaw = redis.call("GET", KEYS[1])
if not metaRaw then return {"not_found"} end
local meta = cjson.decode(metaRaw)
local now = tonumber(ARGV[1])
local claimID = ARGV[3]
local claimUntil = tonumber(ARGV[4])

if meta.status == "consumed" then return {"consumed"} end
if meta.status == "expired" then return {"expired"} end
if meta.status == "deleted" then return {"not_found"} end

if meta.status == "opening" then
	local openingUntil = tonumber(meta.openingUntilUnix) or 0
	if openingUntil > now then
		return {"consumed"}
	end
end

local payloadRaw = redis.call("GET", KEYS[2])
if not payloadRaw then
	meta.status = "expired"
	redis.call("SET", KEYS[1], cjson.encode(meta), "EX", tonumber(ARGV[2]))
	return {"expired"}
end

meta.status = "opening"
meta.openingID = claimID
meta.openingUntilUnix = claimUntil
redis.call("SET", KEYS[1], cjson.encode(meta), "KEEPTTL")

return {"ok", payloadRaw}
\end{minted}